\section{Introduction}\label{sec:intro}
Applications that combine a learned scorer, rules, a small language model (LLM) and a reinforcement-learning (RL) policy are often described as multi-agent systems, although a single controller usually calls each component in a fixed order. Whether the label is earned changes what should be tested. If components are peers, failures arise between agents~\cite{mast}; if one controller coordinates them, the questions are about contracts: who may decide what, who may write which state, what happens when a component misbehaves, and what coordination costs.

We study a spoken technical-interview practice system whose code we audit rather than assume. Before this work its execution path was a hub-orchestrated pipeline with one learned decision-maker (a PPO difficulty policy); the LLM only phrased text and every other decision was a threshold rule in the hub. We \emph{extend} that system, without replacing its workflow, evaluator, PPO policy or frontend/backend contract, so that two specialised, stateful agents take part in real decisions: a Strategy Agent (difficulty proposals and follow-up probes driven by a private weakness ledger) and a Coaching Agent (feedback depth, hints and emphasis driven by a private coaching history). The agents only \emph{propose}. Typed, validated proposals pass a deterministic coordinator and arbitrator; the hub remains the sole writer of session state. The system stays centrally orchestrated; nothing here is decentralised, and the PPO policy is not multi-agent RL.

We ask:
\begin{description}
\item[RQ1 (coordination)] When bounded-authority agents propose decisions that the original workflow takes by rule or by PPO, how often and how do they change the session, and how does arbitration resolve disagreement with PPO and the guardrails?
\item[RQ2 (containment)] Which validation and arbitration mechanisms are necessary to contain agent faults (crash, timeout, malformed, forbidden or stale output, valid-but-harmful proposals, ungrounded plans)?
\item[RQ3 (cost)] What do the agents and the extra checks cost in latency and calls?
\item[RQ4 (carry-over)] How do findings of an earlier audit of this system survive correction (a prediction count, a validation claim, a grounding artefact, scorer and breaker semantics)?
\end{description}

\noindent\textbf{Contributions.} (i)~A design and implementation that adds stateful, bounded-authority agents to an existing orchestrated workflow behind a feature flag, with typed messages, an authority matrix enforced in code (Sect.~\ref{sec:design}--\ref{sec:arbitration}). (ii)~Registered, replay-based experiments comparing the original workflow, the full multi-agent configuration and ablations, a fault-injection study of the safeguards, and a follow-up experiment that isolates the confidence threshold from the other arbitration checks (Sect.~\ref{sec:method}--\ref{sec:results}). (iii)~Corrections of earlier findings with reproducible re-runs (Sect.~\ref{sec:carry}).

\noindent\textbf{What kind of contribution each is.} The agents, typed contracts and authority matrix are \emph{architecture}; arbitration with a confidence threshold and guardrails-last is the \emph{mechanism}; the replay and fault-injection results are \emph{empirical findings about coordination and containment} on constructed data; whether the agents improve interviews or learning is \emph{future work} that needs a live model and human candidates.

\noindent\textbf{Not claimed.} No LLM ran in any experiment; the agents are rule- and ledger-based; all evidence is replayed recorded evaluator output or injected faults on 64 constructed answers. We make no claim about candidate learning, coaching quality, evaluator validity or that agents improve PPO (Sect.~\ref{sec:limits}).
